\documentclass[12pt,a4paper]{report}

\usepackage[
    left=25mm,
    right=10mm,
    top=20mm,
    bottom=20mm
]{geometry}

\usepackage[T2A]{fontenc}
\usepackage[utf8]{inputenc}
\usepackage[russian]{babel}
\usepackage{tempora}
\usepackage{tabularx}
\usepackage{array}
\usepackage{graphicx}
\usepackage{hyperref}
\hypersetup{colorlinks=true, linkcolor=black, citecolor=black, urlcolor=blue}

\linespread{1.3}

\usepackage{fancyhdr}
\pagestyle{fancy}
\fancyhf{}
\fancyfoot[C]{\thepage}
\renewcommand{\headrulewidth}{0pt}

\usepackage{longtable}

\begin{document}

% ============================================================
% ТИТУЛЬНЫЙ ЛИСТ ДНЕВНИКА
% ============================================================
\begin{titlepage}
\begin{center}

\textbf{МИНИСТЕРСТВО ЦИФРОВОГО РАЗВИТИЯ, СВЯЗИ И МАССОВЫХ КОММУНИКАЦИЙ РОССИЙСКОЙ ФЕДЕРАЦИИ}

\vspace{0.3em}

\parbox{0.9\textwidth}{\centering\small
Ордена Трудового Красного Знамени федеральное государственное \\
бюджетное образовательное учреждение \\
высшего образования
}

\vspace{0.3em}

\textbf{«Московский технический университет связи и информатики»} \\
(МТУСИ)

\vfill

\textbf{\Large Д Н Е В Н И К} \\
по производственной практике (проектно-технологической)

\vfill

\begin{flushright}
\begin{minipage}{0.55\textwidth}
\raggedleft

студента \hrulefill \\
\makebox[\linewidth]{\hrulefill} \\
Жуков Никита Александрович \hrulefill \\
\makebox[\linewidth]{\hrulefill} \\
группы БПИ2304 \hrulefill \\
\makebox[\linewidth]{\hrulefill}
\end{minipage}
\end{flushright}

\vfill

\begin{center}
Москва, 2026
\end{center}

\end{center}
\end{titlepage}

% ============================================================
% НАПРАВЛЕНИЕ
% ============================================================
\newpage
\thispagestyle{empty}

\begin{center}
\parbox{0.9\textwidth}{\centering\small
Ордена Трудового Красного Знамени федеральное государственное \\
бюджетное образовательное учреждение \\
высшего образования \\
\textbf{«Московский технический университет связи и информатики»} (МТУСИ)
}
\end{center}

\vspace{1em}

на основании договора \hrulefill

\vspace{1em}

Направляет студента \hrulefill \\
Жуков Никита Александрович \\
\makebox[\linewidth]{\hrulefill} \\
(фамилия, имя, отчество)

группа БПИ2304 \hrulefill

\vspace{1em}

для прохождения \\
производственной практики (проектно-технологической) \hrulefill \\
(вид и тип практики)

на базе ООО «ВИЖНЛАБС» \hrulefill \\
(наименование базы практики)

Срок 02.06.2026–03.07.2026 \hrulefill

\vspace{3em}

\hfill\begin{minipage}{0.5\textwidth}
\raggedleft
Декан факультета ИТ \hrulefill \\
\hspace{2em} (подпись)
\end{minipage}

\vspace{1em}

Печать МТУСИ

\vspace{3em}

Прибыл в организацию «02» июня 2026 г. \\
\hrulefill \\
(подпись)

\vspace{1em}

Выбыл из организации «03» июля 2026 г. \\
\hrulefill \\
(подпись)

\vspace{1em}

Печать (организации)

% ============================================================
% ДНЕВНИК (ТАБЛИЦА)
% ============================================================
\newpage
\thispagestyle{empty}

\begin{center}
\textbf{\Large Д Н Е В Н И К}
\end{center}

\vspace{1em}

\begin{tabularx}{\textwidth}{|p{2.2cm}|X|c|}
\hline
\textbf{Дата} & \textbf{Содержание выполняемых работ} & \textbf{Подпись} \\
\hline
02.06.2026 & Проведение инструктажа и ознакомление с требованиями охраны труда, техники безопасности, пожарной безопасности, а также правилами внутреннего трудового распорядка. Подписание документов и NDA. Ознакомление с программой практики и требованиями к оформлению ее результатов. & \\
\hline
03.06.2026 & Анализ предметной области омниканальных решений. Исследование существующих аналогов (Front, respond.io, Zendesk). Формирование перечня функциональных требований. & \\
\hline
04.06.2026 & Проектирование общей архитектуры приложения на основе Clean Architecture. Определение технологического стека (Python, Litestar, PostgreSQL, React, TypeScript). & \\
\hline
05.06.2026 & Проектирование структуры базы данных PostgreSQL. Определение основных сущностей (messages, senders, clients, curators), их атрибутов и связей. & \\
\hline
08.06.2026 & Настройка среды разработки. Создание структуры серверной части проекта. Настройка pre-commit хуков, линтеров (ruff), системы статической типизации (mypy). & \\
\hline
09.06.2026 & Реализация доменной модели: класс Message, перечисления ChannelType и MessageType. & \\
\hline
10.06.2026 & Реализация портов (интерфейсов) MessageEngine, MessageRepository, DatabaseGateway. & \\
\hline
11.06.2026 & Реализация бизнес-логики: сервисы MessageService, ClientService, AuthService. & \\
\hline
15.06.2026 & Разработка ORM-моделей SQLAlchemy 2.x для таблиц базы данных. Реализация репозиториев. & \\
\hline
16.06.2026 & Реализация DI-контейнера (punq) для связывания компонентов. Настройка точки входа приложения Litestar. & \\
\hline
17.06.2026 & Интеграция с Telegram Bot API. Реализация TelegramEngine: получение сообщений, обработка вложений (фото, документы), отправка ответов. & \\
\hline
18.06.2026 & Интеграция объектного хранилища MinIO. Реализация загрузки медиафайлов из Telegram в S3-хранилище. & \\
\hline
19.06.2026 & Интеграция электронной почты по протоколам IMAP и SMTP. Реализация EmailEngine: получение писем, извлечение MIME-структуры, отправка ответов. & \\
\hline
22.06.2026 & Реализация PollingOrchestrator с асинхронным опросом каналов. Паттерн Observer для уведомлений о новых сообщениях. & \\
\hline
23.06.2026 & Реализация Abstract Factory для создания движков каналов. Регистрация фабрик Telegram и Email. & \\
\hline
24.06.2026 & Разработка REST API: эндпоинты для сообщений, клиентов, кураторов, назначений, аутентификации. WebSocket для уведомлений в реальном времени. & \\
\hline
25.06.2026 & Реализация аутентификации на основе JWT (HS256) и bcrypt. Настройка middleware для проверки токенов. & \\
\hline
26.06.2026 & Разработка клиентского приложения на React 19 + TypeScript + Vite. Создание компонентов интерфейса: левая панель (список контактов) и правая панель (диалог). & \\
\hline
29.06.2026 & Подключение фронтенда к REST API и WebSocket. Реализация отображения сообщений и отправки ответов. & \\
\hline
30.06.2026 & Написание модульных и интеграционных тестов с использованием pytest и pytest-asyncio. & \\
\hline
01.07.2026 & Развёртывание приложения с использованием Docker Compose (PostgreSQL 17, MinIO, серверное приложение). & \\
\hline
02.07.2026 & Комплексное тестирование приложения. Оформление отчётной документации. & \\
\hline
03.07.2026 & Подведение итогов практики. Дифференцированный зачёт. Сдача отчётных документов. & \\
\hline
\end{tabularx}

% ============================================================
% ИНДИВИДУАЛЬНОЕ ЗАДАНИЕ
% ============================================================
\newpage
\thispagestyle{empty}

\begin{center}
\textbf{ИНДИВИДУАЛЬНОЕ ЗАДАНИЕ} \\
по производственной (проектно-технологической) практике \\
(вид и тип практики)
\end{center}

\vspace{1em}

Для студента \hrulefill \\
Жуков Никита Александрович \hrulefill \\
(Ф.И.О.)

Направление подготовки 09.03.04 Программная инженерия \\
Направленность (профиль) подготовки Разработка и сопровождение программного обеспечения

Срок прохождения практики с 02 июня 2026 г. по 03 июля 2026 г.

\vspace{1em}

1. В результате прохождения практики студент должен овладеть следующими компетенциями:
\begin{itemize}
    \item ПК-1 --- Способен разрабатывать требования и проектировать программно-аппаратные комплексы;
    \item ПК-1.1 --- Понимает современные методы и инструментарий проектирования программно-аппаратных комплексов;
    \item ПК-1.2 --- Применяет современные методы и технологии для проектирования, разработки и обеспечения качества разработки программно-аппаратных комплексов для решения задач профессиональной деятельности;
    \item ПК-2 --- Способен осуществлять концептуальное, функциональное и логическое проектирование систем среднего и крупного масштаба и сложности;
    \item ПК-2.1 --- Понимает основы математического аппарата для осуществления проектирования систем разной сложности;
    \item ПК-2.2 --- Применяет теоретические знания и современные технологии для проведения проектирования (логического, функционального) систем разной сложности;
    \item ПК-2.3 --- Разрабатывает и модифицирует архитектуру информационных систем в соответствии с требованиями.
\end{itemize}

\vspace{1em}

2. Индивидуальные задания:

\hrulefill

\parbox{0.95\textwidth}{
Спроектировать и разработать омниканальную платформу для центра обучения, обеспечивающую приём и обработку обращений из различных каналов связи в едином интерфейсе куратора. Реализовать серверную часть на Python (Litestar) с архитектурой Clean Architecture, клиентский интерфейс на React 19 с TypeScript, интеграцию каналов Telegram и Email. Выполнить развёртывание с использованием Docker Compose.
}

\vspace{3em}

\begin{tabularx}{\textwidth}{p{5cm} X p{3cm}}
Руководитель практики от кафедры МТУСИ &
\makebox[5cm]{\hrulefill} Старший преподаватель кафедры МКиИТ, Манжос М.С. &
02.06.2026 \\
& (должность, Ф.И.О., подпись) & (дата)
\end{tabularx}

\vspace{1em}

\begin{tabularx}{\textwidth}{p{5cm} X p{3cm}}
Согласовано: &
\makebox[5cm]{\hrulefill} & \\
Руководитель практики от организации &
Жариков Александр Романович &
02.06.2026 \\
& Руководитель отдела аналитики и архитектуры ООО «ВИЖНЛАБС» & (дата)
\end{tabularx}

\vspace{1em}

\begin{tabularx}{\textwidth}{p{5cm} X p{3cm}}
Студент &
\makebox[5cm]{\hrulefill} &
02.06.2026 \\
& Жуков Никита Александрович & (дата) \\
& (Ф.И.О., подпись) &
\end{tabularx}

% ============================================================
% ПЛАН (РАБОЧИЙ ГРАФИК)
% ============================================================
\newpage
\thispagestyle{empty}

\begin{center}
\textbf{ПЛАН (рабочий график)} \\
производственной (проектно-технологической) практики
\end{center}

\vspace{1em}

Для студента \hrulefill \\
Жуков Никита Александрович \hrulefill \\
(Ф.И.О.)

Направление подготовки 09.03.04 Программная инженерия \\
Направленность (профиль) Разработка и сопровождение программного обеспечения

Срок прохождения практики с 02 июня 2026 г. по 03 июля 2026 г.

в ООО «ВИЖНЛАБС» \hrulefill \\
(наименование организации)

\vspace{1em}

\begin{tabularx}{\textwidth}{|p{2.8cm}|X|p{3cm}|}
\hline
\textbf{Период (дата)} & \textbf{Содержание практики (наименование работ)} & \textbf{Планируемые результаты (освоенные компетенции)} \\
\hline
02.06 & Проведение инструктажа. Подписание документов и NDA. Ознакомление с программой практики. & ПК-1, ПК-1.1 \\
\hline
02.06 & Получение задания на практику. Формирование технического задания. & ПК-1, ПК-2.1 \\
\hline
03.06–05.06 & Анализ предметной области омниканальных решений. Изучение аналогов. Проектирование архитектуры и базы данных. & ПК-1.2, ПК-2.3 \\
\hline
08.06–14.06 & Работа над индивидуальным заданием: реализация доменной модели, портов, сервисов и DI-контейнера. & ПК-1.1, ПК-2.2 \\
\hline
15.06–21.06 & Работа над индивидуальным заданием: интеграция каналов Telegram и Email, реализация адаптеров, фабрики движков. & ПК-1.2, ПК-2.3 \\
\hline
22.06–28.06 & Работа над индивидуальным заданием: разработка REST API, WebSocket, клиентского приложения на React. & ПК-1.1, ПК-1.2 \\
\hline
29.06–02.07 & Тестирование, развёртывание с Docker Compose, оформление отчётной документации. & ПК-1.2, ПК-2.2 \\
\hline
03.07 & Дифференцированный зачёт. Сдача отчётных документов. & ПК-1, ПК-2.3 \\
\hline
\end{tabularx}

\vspace{2em}

\begin{tabularx}{\textwidth}{p{5cm} X p{3cm}}
Руководитель практики от кафедры МТУСИ &
\makebox[5cm]{\hrulefill} Старший преподаватель кафедры МКиИТ, Манжос М.С. &
02.06.2026 \\
& (должность, Ф.И.О., подпись) & (дата)
\end{tabularx}

\vspace{1em}

\begin{tabularx}{\textwidth}{p{5cm} X p{3cm}}
Согласовано: &
\makebox[5cm]{\hrulefill} & \\
Руководитель практики от организации &
Жариков Александр Романович &
02.06.2026 \\
& Руководитель отдела аналитики и архитектуры ООО «ВИЖНЛАБС» & (дата)
\end{tabularx}

\vspace{1em}

\begin{tabularx}{\textwidth}{p{5cm} X p{3cm}}
Студент &
\makebox[5cm]{\hrulefill} &
02.06.2026 \\
& Жуков Никита Александрович & (дата) \\
& (Ф.И.О., подпись) &
\end{tabularx}

% ============================================================
% ОТЗЫВ
% ============================================================
\newpage
\thispagestyle{empty}

\begin{center}
\textbf{ОТЗЫВ} \\
о прохождении производственной (проектно-технологической) практики
\end{center}

\vspace{1em}

Студент \hrulefill \\
Жуков Никита Александрович \hrulefill \\
(Ф.И.О.)

Направление подготовки 09.03.04 Программная инженерия \\
Направленность (профиль) Разработка и сопровождение программного обеспечения

Срок прохождения практики с 02 июня 2026 г. по 03 июля 2026 г.

в ООО «ВИЖНЛАБС» \hrulefill \\
(наименование организации/структурного подразделения)

\vspace{1em}

\textbf{Отзыв о прохождении студентом практики} (выполнение индивидуального задания, плана практики, отношение к работе, трудовая дисциплина, овладение производственными навыками, участие в научно-исследовательской, рационализаторской, общественной работе и др.).

\vspace{1em}

\parbox{0.95\textwidth}{
Во время прохождения производственной практики в ООО «ВИЖНЛАБС» студент группы БПИ2304 Жуков Никита Александрович выполнил все задания, предусмотренные индивидуальным планом, в установленные сроки и в полном объёме.

За время практики зарекомендовал себя как ответственный, внимательный и пунктуальный сотрудник. Проявил инициативность и заинтересованность в работе, успешно справлялся с поставленными задачами.

Продемонстрировал хорошие профессиональные навыки, умение грамотно анализировать поставленные задачи и находить эффективные решения. При выполнении заданий проявлял самостоятельность, организованность и творческий подход. Умеет работать в коллективе, пользуется уважением и доверием коллег.

В целом, показал себя как компетентный и перспективный специалист, готовый к профессиональной деятельности по специальности.
}

\vspace{2em}

\hrulefill

\vspace{1em}

(печать организации)

\vspace{3em}

\textbf{Оценка результатов прохождения практики и выполнения индивидуального задания} \\
\hrulefill \\
(отлично, хорошо, удовлетворительно, неудовлетворительно)

\vspace{1em}

Ф.И.О., подпись \hrulefill \\
Жариков Александр Романович \hrulefill

\vspace{1em}

«\hrulefill » \hrulefill 2026 г.

\end{document}
